\scanpage{34}
On trouve ensuite

$\widetilde S=A^{\vcenter{\hbox{\tikz[x=.15cm,y=.15cm]{\draw[->] (0,0)--(0,1.2);\draw[->] (0,0)--(1.3,0);}}}}/\rho_s\simeq$ ens des couples $(s,\omega)$, $s\in S$, $\omega$ une orientation du \deleted{\unsure{lien}}
polygône $\mathcal X_{\geq s}$ des $a\in A$ et $f\in F$ incidents à $s$
[i.e. un \inserted{des 2} ordres circulaires sur l'ens des $a\in A$ incidents à $s$, qui définissent sa structure polygonale]

$\widetilde F=A^{\vcenter{\hbox{\tikz[x=.15cm,y=.15cm]{\draw[->] (0,0)--(0,1.2);\draw[->] (0,0)--(1.3,0);}}}}/\rho_f=$ ens des couples $(f,\unsure{\varpi|\bar\omega})$, $f\in F$, $\varpi$ une orientation du
polygône $\mathcal X_{\leq f}$ des \deleted{$\unsure{f}$} $a\in A$ et $s\in S$ incidents à $f$,
[i.e. un des 2 ordres circ. sur l'ens. des
$a\in A$ incidents à $f$, qui définissent sa structure polygonale]

$\widetilde A=A^{\vcenter{\hbox{\tikz[x=.15cm,y=.15cm]{\draw[->] (0,0)--(0,1.2);\draw[->] (0,0)--(1.3,0);}}}}/\sigma=$ ens des couples $(a,u)$, $a\in A$, $u\in\vec A(a)\wedge A^\uparrow(a)$,
où $\vec A(a)$ ($A^\uparrow(a)$) est l'ens des \unsure{axes|arcs|sens} tang.
(axes transverses) portés par l'arête $a$.

\[
A^\wedge=\mathrm{Im}\,(A^{\vcenter{\hbox{\tikz[x=.15cm,y=.15cm]{\draw[->] (0,0)--(0,1.2);\draw[->] (0,0)--(1.3,0);}}}}\to S\times F)
=\begin{aligned}[t]&\text{rel. d'incidence}\\&\text{entre }S\text{ et }F\end{aligned}
\]

\diagram{\begin{tikzcd}[ampersand replacement=\&,row sep=large,column sep=small]
\&\& A^{\vcenter{\hbox{\tikz[x=.15cm,y=.15cm]{\draw[->] (0,0)--(0,1.2);\draw[->] (0,0)--(1.3,0);}}}} \arrow[dll,"\rho_s"']\arrow[dl,"\unsure{\sigma_2|\sigma}"']\arrow[d,"\sigma" description]\arrow[dr,"\sigma_1"]\arrow[drr,"\rho_f"]\&\&\\
\widetilde S\arrow[d,"\sigma_1"']\&\vec A\arrow[dl]\arrow[dr,"\sigma_1"']\&
\widetilde A\arrow[d,"\sigma_1"',"\sigma_2"]\&A^\uparrow\arrow[dl,"\sigma_2"]\arrow[dr]\&\widetilde F\arrow[d,"\sigma_2"]\\
S\&\&A\&\&F
\end{tikzcd}}

\[
\begin{aligned}
A^{\vcenter{\hbox{\tikz[x=.15cm,y=.15cm]{\draw[->] (0,0)--(0,1.2);\draw[->] (0,0)--(1.3,0);}}}}&\simeq\widetilde S\times_S\vec A\simeq\deleted{\unsure{\widetilde F}\times_F}A^\uparrow\times_F\widetilde F\\
&\simeq\vec A\times_A A^\uparrow\\
&\simeq\vec A\times_A\widetilde A\simeq\widetilde A\times_A A^\uparrow\\
&\simeq\widetilde S\times_S A^\wedge\simeq A^\wedge\times_F\widetilde F
\end{aligned}
\]

\diagram{\begin{tikzcd}[ampersand replacement=\&,row sep=small,column sep=small]
\&A^{\vcenter{\hbox{\tikz[x=.15cm,y=.15cm]{\draw[->] (0,0)--(0,1.2);\draw[->] (0,0)--(1.3,0);}}}}\arrow[dl]\arrow[d]\arrow[dr,dash]\&\\
\widetilde S\arrow[d]\&A^\wedge\arrow[dl]\arrow[dr]\&\widetilde F\arrow[d,dash]\\
S\&\&F
\end{tikzcd}}

$A^{\vcenter{\hbox{\tikz[x=.15cm,y=.15cm]{\draw[->] (0,0)--(0,1.2);\draw[->] (0,0)--(1.3,0);}}}}$ s'identifie aux rel. d'incidences
\[
\begin{aligned}
&\text{entre }(S,A,F)\quad\text{simultanément}\\
&\text{entre }\widetilde S,\widetilde A,\\
&\text{entre }\widetilde A,\widetilde F,\\
&\text{entre }\widetilde S,\widetilde F,
\end{aligned}
\]
% archive sheet number: 33
